\section{Evaluation Design}\label{sec:evaluation}

This section specifies the replacement evaluation campaign. The
measurements and application traces will be generated anew; figures in
Motivation illustrate mechanisms and do not supply campaign observations.
The campaign tests workflow-level protection, the incremental value of
Peer-HBM after DRAM coverage, and healthy serving costs.

Figure~\ref{fig:evaluation-protocol} separates application capture and
quality validation from the paired serving comparison.

\begin{figure*}[t]
\centering
\includegraphics[width=168mm]{figures/evaluation-protocol.pdf}
\caption{New evaluation protocol. Real application execution supplies
full Agent histories and explicit dependencies. Quality is checked
independently. Serving replay preserves dependency gates and compares
healthy F0 and faulted F1 runs at the same inputs, arrivals, and resource
budget, measuring paired completion, protection cost, and interference.}
\label{fig:evaluation-protocol}
\Description{Public tasks and frozen materials feed real model execution.
Its outputs feed independent application quality checks and a captured
trace. Dependency-gated replay runs healthy and worker-fault branches,
which feed paired job-completion and cost analysis.}
\end{figure*}

\subsection{Applications and Trace Collection}\label{sec:eval-workloads}

We instantiate the research-report and repository-review Agent-task DAGs
introduced in Section~\ref{sec:background:serving} with independent public tasks. Research inputs come from DeepResearchBench
and compatible tasks in DeepResearch Bench II
\cite{du2025deepresearch,li2026deepresearchii}. Eligibility is defined by
the need to compare alternatives, application criteria, and contextual
constraints, before observing recovery outcomes. The report application
is not limited to software markets. Repository inputs use SWE-bench
Verified issues and their pre-fix code versions \cite{swebenchverified};
the application produces review and change plans, rather than executing
the original patch-resolution benchmark.

Each Agent obtains all predecessor outputs before starting and retains its
complete local message history. Material groups, analysis, checking, and
necessary revision produce separate invocations. The main application
variant uses frozen source materials without tools. Preparation time is
reported separately; local-file and online-retrieval variants retain their
actual tool costs. Context overflow, quality failures, and excluded tasks
remain in the accounting.

The collector records job, Agent, invocation and predecessor identities;
full messages and serialized input/output tokens; reusable token-prefix
lengths; submit, admission, GPU-start and completion timestamps; actual
worker/TP placement; tools; checkpoint frontiers; native selections;
transfer timelines; and failure/retry events. Real model generation
establishes application behavior and quality. Deterministic serving replay
fixes those captured token sequences while preserving Agent-local order,
predecessor completion gates, and external job arrivals. It does not fix
F1 successor submissions to F0 absolute times.

\subsection{Platform and Baselines}\label{sec:eval-setup}

The primary platform requires at least two compatible replicas with a
surviving peer in the stated protection domain. We freeze actual GPU,
host memory, PCIe/NVLink/network topology, TP layout, precision, model and
tokenizer hashes, native kernels, scheduling parameters, and failure
detection. Native-sparse candidates must pass generation-quality and
recovery-compatibility pilots; a second supported model tests generality.
Dense-model tests, if included, evaluate Agent protection separately from
sparse-region claims.

All policies share normal caching, batching, durability, compatible
routing, and physical byte budgets. The core baselines are restart without
extra replicas; fixed-interval budgeted checkpoints without Slack;
reconstruction-cost priority; \system with DRAM only; \system with fixed
HBM quota; and the full adaptive design. Full DRAM protection is a feasible
reference only where it fits; an unlimited-capacity reference is labeled
separately. Available implementations of
D\'ej\`aVu, LUMEN, and GhostServe are external comparisons when faithfully
adapted \cite{strati2024dejavu,cao2026lumen,jayakody2026ghostserve};
an internal inspired policy is not a reproduction.

Normal workflow schedulers and caches are complementary
\cite{lin2024parrot,pan2025kvflow,ni2026topas}. A recovery comparison on
such a backend keeps that backend fixed across policies. Recovery assets
already reachable through normal caching count toward the baseline.

\subsection{Paired Failure and End-to-End Performance}\label{sec:eval-main}

We calibrate low, medium, and high arrival rates on disjoint healthy tasks,
then save identical arrival sequences across policies. Every policy runs
both without failure (F0) and with failure (F1). A controlled failure is
bound to a job, Agent, invocation, phase, and accepted-token or prefill
landmark. Early, middle, and late failures, mid-prefill and decode, and
one- or multiple-Agent victim sets are separate cells. Natural wall-time
worker failures additionally test deployment behavior, recording the actual
victim set when policy-dependent schedules differ.

For policy $p$, the paired penalty is
$P_p=W_{p,\mathrm{F1}}-W_{p,\mathrm{F0}}$.
We report absolute F1 WCT, differences in $P_p$, and F0 overhead. Primary
metrics include mean, median, and supported tail WCT; completion rate;
job goodput; TTFT and TPOT; actual DRAM/HBM bytes; and copying and controller
cost. A worker failure can affect other jobs through takeover load and
shared transfer queues. We therefore report directly interrupted,
indirectly affected, and all submitted jobs, with failure-relative latency
and backlog timelines. Unfinished jobs are not silently removed.

\subsection{Mechanism and Component Experiments}\label{sec:eval-native}

Workload characterization reports calls per Agent, token-prefix growth,
actual execution overlap, inter-call gaps, and measured tool share.
Cost experiments compare reconstruction, full DRAM, and DRAM plus HBM on
the same actual prefixes, accepted tails, retry target, and load. They
include auxiliary state and all TP shards.

To isolate workflow value, we compare reconstruction-cost priority,
zero-Slack protection, causal estimated Slack, and an explicitly offline
oracle at matched budgets. The joint distribution of failure-time Slack
and reconstruction delay tests the premise rather than substituting
branch-duration gaps.

To isolate HBM value, we hold protected Agents and DRAM manifests fixed.
Causal native-count, recency, and random selectors receive equal physical
bytes, layer ranges, and mandatory local/sink allocations. A future-demand
oracle supplies an offline reference. We count distinct required
layer/block/shard bytes and measure actual exposed H2D waiting, misses,
tail replay, and end-to-end retry latency. Query frequencies or potential
coverage alone do not establish recovery savings.

An eight-cell ablation varies workflow selection, fast-region placement,
and quota adaptation. Dynamic-load tests compare adaptive, conservative
fixed, and aggressive fixed quotas, reporting safe capacity, target quota,
actual resident replicas, churn, and foreground latency on the same time
axis. Prediction experiments use task/source-disjoint training and held-out
causal features, reporting remaining-duration errors, Slack bias,
missed-protection rates, WCT regret, and controller overhead.

\subsection{Limits, Correctness, and Statistical Protocol}\label{sec:eval-limits}

Sensitivity tests vary protection threshold $\lambda$, checkpoint chunk
$Z$, block count $K$, depth $L$, capacity, measured effective bandwidth,
concurrency, and worker count. Fixed $K\times L$ and equal-byte comparisons
are distinguished. Short histories, sufficient DRAM bandwidth, scattered
accesses, and poor prediction are retained even when HBM or protection
has no benefit.

Fault tests cover pending copies, incomplete TP shards, quota shrink,
unavailable peers, and repeated failures. Recovery uses the last complete
committed frontier; the replay tail can exceed $Z$ when copying lags.
Token equivalence, durable accepted outputs, and tool idempotence are
checked separately from timing. The failure-domain claim matches the
actual injection: process termination is not labeled host loss.

A pilot determines sample size and test precision before tuning on the
held-out tasks. Repeated calls, failure positions, and timing repetitions
are not independent application samples. We use paired estimates and
task/source- and arrival-block-aware confidence intervals, publish
denominators and failures, and avoid unsupported p99 claims on small
cohorts. Research quality uses held-out evidence/rubric checks; repository
review uses independently verified causes, source references, and
actionable changes/tests. Reference reports, gold patches and answer tests
are excluded from generation inputs.

Healthy cost and fault benefit are computed on the same population and
objective. If $H=W_{\mathrm{A},0}-W_{\mathrm{B},0}>0$ and
$D=W_{\mathrm{B},1}-W_{\mathrm{A},1}>0$, a specified binary mixture has
break-even probability $H/(H+D)$. This conditional analysis does not
estimate production failure frequency. The artifact release binds every
plot to its raw trace hashes, model configuration, paired fault schedule,
statistical unit, and physical or analytical evidence class.
